%% ============================================================================
%% §7.2 ASIC Verification  --  Willie
%% ============================================================================
\subsection{ASIC Verification: Three-Tier + SDF Signoff}
\label{sec:verify:asic}

The ASIC path is verified at four progressively stricter levels; results are summarised in Table~\ref{tab:asic_verify_results}.

\begin{table}[H]
\centering
\caption{ASIC Verification Results}
\label{tab:asic_verify_results}
\begin{tabular}{clcc}
\toprule
\textbf{Level} & \textbf{Stimulus} & \textbf{Delay model} & \textbf{Result} \\
\midrule
L1 & RTL                & zero-delay & 50 / 50 MATCH \\
L2 & Genus syn netlist  & zero-delay & 50 / 50 MATCH \\
L3 & Innovus post-route & zero-delay & 50 / 50 MATCH \\
L4 & Innovus post-route & SDF annotated & 50 / 50 MATCH, 0 violations \\
\midrule
\multicolumn{4}{l}{End-to-end embedding (RTL, 13 face images): 128 / 128 exact, MaxDiff = 0} \\
\bottomrule
\end{tabular}
\end{table}

Levels~L1--L3 collectively establish that the design behaves identically at the RTL, synthesised, and post-route stages. A formal Logic Equivalence Checker (Conformal LEC, Formality) would normally produce this result in seconds; without one available we substitute a full 50-layer gate-level simulation against the C-model golden hex vectors. The trade-off is wall-clock time. Table~\ref{tab:gl_runtime} summarises how long each of the three gate-level variants took on the same machine.

\begin{table}[H]
\centering
\caption{Gate-Level Simulation Wall-Clock Comparison}
\label{tab:gl_runtime}
\begin{tabular}{lllcr}
\toprule
\textbf{Run} & \textbf{Netlist} & \textbf{Delay model} & \textbf{Layers done} & \textbf{Wall-clock} \\
\midrule
\texttt{make gl}      & Genus syn         & zero-delay      & 50 / 50 & $\sim$78\,h \\
\texttt{make pnr}     & Innovus post-route & zero-delay     & 50 / 50 & $\sim$84\,h \\
\texttt{make pnr\_sdf} & Innovus post-route & SDF annotated & 50 / 50 & $\sim$163\,h \\
\bottomrule
\end{tabular}
\end{table}

Post-route simulation runs about $1.1\times$ slower than synthesised because of the added CTS buffers, filler cells, and hold-fix buffers in the post-route netlist. SDF backannotation adds another $1.8\times$ on top: every net carries an RC delay and every flop fires four timing checks (setup, hold, recovery, removal). High-resolution early layers (0--2) dominate all three runs and account for roughly half the wall-clock time.

L4 uses the Innovus-exported SDF (\texttt{mfn\_frontend\_top\_v3\_final.sdf}, 117\,MB) backannotated into the same gate-level simulation. After 163\,h of wall-clock time the full 50-layer run completed with all 50 layers bit-exact against the C-model and zero \texttt{TCHK\_VIOL}, setup-violation, or hold-violation events. In the absence of PrimeTime SI this is the strongest signoff we can produce: the netlist runs MobileFaceNet correctly at 200\,MHz with real cell and wire delays and live setup/hold checks.

The verification flow itself had one early failure worth recording. The first gate-level run reported 50/50 mismatches even though RTL was 50/50 MATCH. A SystemVerilog \texttt{bind}-based probe traced the cause to a testbench race: \texttt{pixel\_in} was driven on \texttt{@(posedge clk)} from the DUT's \texttt{x\_out} output, but the gate-level netlist has delta-cycle delay on port outputs, so the testbench was latching the previous-cycle address. Moving the drive to \texttt{@(negedge clk)} (after \texttt{x\_out} has settled) restored all 50 layers. The same testbench is now used for every gate-level and SDF run.

Limitations to keep in mind: this is a single typical-corner signoff (FreePDK45 does not ship multi-corner Liberty), there is no signal-integrity or cross-talk analysis, and the equivalence is shown by simulation rather than formal proof.
